\documentclass[border=3pt,tikz]{standalone}
\usepackage{amsmath} % for \dfrac
\usepackage{physics}
\usepackage{tikz,pgfplots}
\usepackage{circuitikz}
\pgfplotsset{compat=1.18}
\usepackage{xcolor}



\begin{document}
\begin{figure}[!ht]
    \centering
    \resizebox{1\textwidth}{!}{%
        \begin{circuitikz}
            \tikzstyle{every node}=[font=\LARGE]
            \draw (5.5,10.5) to[C] (5.5,9.25);
            \draw[, line width=1.2pt] (7.5,10.25) to[short] (6.25,10.25);
            \draw (7.75,10.25) to[short] (6.25,10.25);
            \draw [, line width=1.2pt](8.75,10.25) to[short] (10,10.25);
            \draw [, line width=1.2pt](10,10.25) to[short] (10,7.75);
            \draw [, line width=1.2pt](6.25,10.25) to[short] (6.25,7.75);
            \draw [line width=1.2pt](7.5,10.25) to[C,l={ \large $C$}] (8.75,10.25);
            \draw [line width=1.2pt](10,7.75) to (10.5,7.75) node[ground]{};
            \draw (7.5,7.75) to[battery1,l=$\Delta V$] (8.75,7.75);
            \draw (8.25,7.5) to[battery1] (8,7.5);
            \draw (6.25,7.75) to[short, -*] (6.25,7.75);
            \draw (6.25,10.25) to[short, -*] (6.25,10.25);
            \draw (10,10.25) to[short, -*] (10,10.25);
            \draw (10,7.75) to[short, -*] (10,7.75);
            \draw [, line width=1.2pt](6.25,7.75) to[short] (7.5,7.75);
            \draw [, line width=1.2pt](8.75,7.75) to[short] (10,7.75);
            \node [font=\large, color={rgb,255:red,255; green,0; blue,0}] at (7.75,7.5) {+};
            \node [font=\large, color={rgb,255:red,0; green,8; blue,255}] at (8.5,7.5) {-};
        \end{circuitikz}
    }%
    \label{fig:my_label}
\end{figure}
\end{document}